\documentclass{amsart}
% For dealing with references we use the comment environment
\usepackage{verbatim}
\newenvironment{reference}{\comment}{\endcomment}
%\newenvironment{reference}{}{}
\newenvironment{slogan}{\comment}{\endcomment}
\newenvironment{history}{\comment}{\endcomment}

% For commutative diagrams we use Xy-pic
\usepackage[all]{xy}

% We use 2cell for 2-commutative diagrams.
\xyoption{2cell}
\UseAllTwocells

% We use multicol for the list of chapters between chapters
\usepackage{multicol}

% This is generally recommended for better output
\usepackage{lmodern}
\usepackage[T1]{fontenc}

% For cross-file-references

% Package for hypertext links:
\usepackage{hyperref}

% For any local file, say "hello.tex" you want to link to please

% Theorem environments.
%
\theoremstyle{plain}
\newtheorem{theorem}[subsection]{Theorem}
\newtheorem{proposition}[subsection]{Proposition}
\newtheorem{lemma}[subsection]{Lemma}

\theoremstyle{definition}
\newtheorem{definition}[subsection]{Definition}
\newtheorem{example}[subsection]{Example}
\newtheorem{exercise}[subsection]{Exercise}
\newtheorem{situation}[subsection]{Situation}

\theoremstyle{remark}
\newtheorem{remark}[subsection]{Remark}
\newtheorem{remarks}[subsection]{Remarks}

\numberwithin{equation}{subsection}

% Macros
%
\def\lim{\mathop{\mathrm{lim}}\nolimits}
\def\colim{\mathop{\mathrm{colim}}\nolimits}
\def\Spec{\mathop{\mathrm{Spec}}}
\def\Hom{\mathop{\mathrm{Hom}}\nolimits}
\def\Ext{\mathop{\mathrm{Ext}}\nolimits}
\def\SheafHom{\mathop{\mathcal{H}\!\mathit{om}}\nolimits}
\def\SheafExt{\mathop{\mathcal{E}\!\mathit{xt}}\nolimits}
\def\Sch{\mathit{Sch}}
\def\Mor{\mathop{\mathrm{Mor}}\nolimits}
\def\Ob{\mathop{\mathrm{Ob}}\nolimits}
\def\Sh{\mathop{\mathit{Sh}}\nolimits}
\def\NL{\mathop{N\!L}\nolimits}
\def\CH{\mathop{\mathrm{CH}}\nolimits}
\def\proetale{{pro\text{-}\acute{e}tale}}
\def\etale{{\acute{e}tale}}
\def\QCoh{\mathit{QCoh}}
\def\Ker{\mathop{\mathrm{Ker}}}
\def\Im{\mathop{\mathrm{Im}}}
\def\Coker{\mathop{\mathrm{Coker}}}
\def\Coim{\mathop{\mathrm{Coim}}}

% Boxtimes
%
\DeclareMathSymbol{\boxtimes}{\mathbin}{AMSa}{"02}

%
% Macros for moduli stacks/spaces
%
\def\QCohstack{\mathcal{QC}\!\mathit{oh}}
\def\Cohstack{\mathcal{C}\!\mathit{oh}}
\def\Spacesstack{\mathcal{S}\!\mathit{paces}}
\def\Quotfunctor{\mathrm{Quot}}
\def\Hilbfunctor{\mathrm{Hilb}}
\def\Curvesstack{\mathcal{C}\!\mathit{urves}}
\def\Polarizedstack{\mathcal{P}\!\mathit{olarized}}
\def\Complexesstack{\mathcal{C}\!\mathit{omplexes}}
% \Pic is the operator that assigns to X its picard group, usage \Pic(X)
% \Picardstack_{X/B} denotes the Picard stack of X over B
% \Picardfunctor_{X/B} denotes the Picard functor of X over B
\def\Pic{\mathop{\mathrm{Pic}}\nolimits}
\def\Picardstack{\mathcal{P}\!\mathit{ic}}
\def\Picardfunctor{\mathrm{Pic}}
\def\Deformationcategory{\mathcal{D}\!\mathit{ef}}

\setlength{\emergencystretch}{3em}
\begin{document}
\title[Stacks Project, Tag 0BQZ]{257 --- Dualizing Ext for a Koszul-regular immersion and its normal determinant}
\maketitle
\noindent Copyright (C) 2005--2025 Johan de Jong. Stacks Project authors. Source: \href{https://stacks.math.columbia.edu/tag/0BQZ}{Tag 0BQZ}.

\noindent Editorial difficulty note (not author text). Difficulty H2: The proof explicitly identifies the dual Koszul complex with a shifted Koszul complex using contraction derivations. It tracks the odd/even sign discrepancy and the top cohomology map to the normal determinant, thereby proving both unique-degree vanishing and canonical identification. This is a structural derived calculation with compatibility obligations..

\noindent Editorial provenance note (not author text). History: excerpt from revision \texttt{a0444\allowbreak 6e57e\allowbreak c1fbc\allowbreak 252a8\allowbreak 71afc\allowbreak ec775\allowbreak 2fb28\allowbreak 07b14}, duality.tex, lines 3543--3632. Reference presentation was adapted; original-author context and core support blocks were retained or added. The mathematical wording is unchanged. Licensed under GNU FDL 1.2 or later; no invariant sections or cover texts. See ../licenses/COPYING. Context blocks reproduce author definitions and supporting results; they are not additional counted problems.

\begin{lemma}
\label{lemma-regular-ideal-ext}
Let $Y$ be a ringed space. Let $\mathcal{I} \subset \mathcal{O}_Y$
be a sheaf of ideals. Set $\mathcal{O}_X = \mathcal{O}_Y/\mathcal{I}$.
If $\mathcal{I}$ is Koszul-regular
(Divisors, Definition \ref{divisors-definition-regular-ideal-sheaf})
then composition on $R\SheafHom_{\mathcal{O}_Y}(\mathcal{O}_X, \mathcal{O}_X)$
defines isomorphisms
$$
\wedge^i(\SheafExt^1_{\mathcal{O}_Y}(\mathcal{O}_X, \mathcal{O}_X))
\longrightarrow
\SheafExt^i_{\mathcal{O}_Y}(\mathcal{O}_X, \mathcal{O}_X)
$$
for all $i$.
\end{lemma}

\begin{definition}
\label{divisors-definition-regular-ideal-sheaf}
\begin{reference}
The concept of a Koszul-regular ideal sheaf was introduced in
\href{https://stacks.math.columbia.edu/bibliography}{Bibliography: SGA6 (Expose VII, Definition 1.4)} where it was called a
regular ideal sheaf.
\end{reference}
Let $X$ be a ringed space. Let $\mathcal{J} \subset \mathcal{O}_X$
be a sheaf of ideals.
\begin{enumerate}
\item We say $\mathcal{J}$ is {\it regular} if for every
$x \in \text{Supp}(\mathcal{O}_X/\mathcal{J})$ there exists an open
neighbourhood $x \in U \subset X$ and a regular sequence
$f_1, \ldots, f_r \in \mathcal{O}_X(U)$ such that $\mathcal{J}|_U$
is generated by $f_1, \ldots, f_r$.
\item We say $\mathcal{J}$ is {\it Koszul-regular} if for every
$x \in \text{Supp}(\mathcal{O}_X/\mathcal{J})$ there exists an open
neighbourhood $x \in U \subset X$ and a Koszul-regular sequence
$f_1, \ldots, f_r \in \mathcal{O}_X(U)$ such that $\mathcal{J}|_U$
is generated by $f_1, \ldots, f_r$.
\item We say $\mathcal{J}$ is {\it $H_1$-regular} if for every
$x \in \text{Supp}(\mathcal{O}_X/\mathcal{J})$ there exists an open
neighbourhood $x \in U \subset X$ and a $H_1$-regular sequence
$f_1, \ldots, f_r \in \mathcal{O}_X(U)$ such that $\mathcal{J}|_U$
is generated by $f_1, \ldots, f_r$.
\item We say $\mathcal{J}$ is {\it quasi-regular} if for every
$x \in \text{Supp}(\mathcal{O}_X/\mathcal{J})$ there exists an open
neighbourhood $x \in U \subset X$ and a quasi-regular sequence
$f_1, \ldots, f_r \in \mathcal{O}_X(U)$ such that $\mathcal{J}|_U$
is generated by $f_1, \ldots, f_r$.
\end{enumerate}
\end{definition}

\begin{lemma}
\label{lemma-regular-immersion-ext}
Let $Y$ be a ringed space. Let $\mathcal{I} \subset \mathcal{O}_Y$
be a sheaf of ideals. Set $\mathcal{O}_X = \mathcal{O}_Y/\mathcal{I}$ and
$\mathcal{N} =
\SheafHom_{\mathcal{O}_Y}(\mathcal{I}/\mathcal{I}^2, \mathcal{O}_X)$.
If $\mathcal{I}$ is Koszul-regular
(Divisors, Definition \ref{divisors-definition-regular-ideal-sheaf}) then
$$
R\SheafHom_{\mathcal{O}_Y}(\mathcal{O}_X, \mathcal{O}_Y) =
\wedge^r \mathcal{N}[-r]
$$
where $r : Y \to \{1, 2, 3, \ldots \}$ sends $y$ to
the minimal number of generators of $\mathcal{I}$ needed in a neighbourhood
of $y$.
\end{lemma}

\begin{proof}
We can use Lemmas \href{https://stacks.math.columbia.edu/tag/0BQW}{lemma ext (Tag 0BQW)} and \ref{lemma-regular-ideal-ext}
to see that we have isomorphisms
$\wedge^i\mathcal{N} \to
\SheafExt^i_{\mathcal{O}_Y}(\mathcal{O}_X, \mathcal{O}_X)$
for $i \geq 0$. Thus it suffices to show that the map
$\mathcal{O}_Y \to \mathcal{O}_X$ induces an isomorphism
$$
\SheafExt^r_{\mathcal{O}_Y}(\mathcal{O}_X, \mathcal{O}_Y)
\longrightarrow
\SheafExt^r_{\mathcal{O}_Y}(\mathcal{O}_X, \mathcal{O}_X)
$$
and that
$\SheafExt^i_{\mathcal{O}_Y}(\mathcal{O}_X, \mathcal{O}_Y)$
is zero for $i \not = r$. These statements are local on $Y$. Thus
we may assume that we have global sections $f_1, \ldots, f_r$ of
$\mathcal{O}_Y$ which generate $\mathcal{I}$ and which form a
Koszul regular sequence. Let $\mathcal{A}^\bullet$
be the Koszul complex on $f_1, \ldots, f_r$ as introduced in the proof of
Lemma \ref{lemma-regular-ideal-ext}. Then
$$
R\SheafHom_{\mathcal{O}_Y}(\mathcal{O}_X, \mathcal{O}_Y) =
\SheafHom^\bullet(\mathcal{A}^\bullet, \mathcal{O}_Y)
$$
by Cohomology, Lemma \href{https://stacks.math.columbia.edu/tag/08DM}{lemma Rhom strictly perfect (Tag 08DM)}.
Denote $1 : \mathcal{A}^\bullet \to \mathcal{O}_Y$ the map of differential
graded $\mathcal{O}_Y$-algebras given by the identity map of
$\mathcal{A}^0 = \mathcal{O}_Y \to \mathcal{O}_Y$ in degree $0$.
With $\delta_j$ as in the proof of Lemma \ref{lemma-regular-ideal-ext}
we get an isomorphism of graded $\mathcal{O}_Y$-modules
$$
\mathcal{O}_Y\langle \delta_1, \ldots, \delta_r\rangle
\longrightarrow
\SheafHom^\bullet(\mathcal{A}^\bullet, \mathcal{O}_Y)
$$
by mapping $\delta_{j_1} \ldots \delta_{j_i}$ to
$1 \circ \delta_{j_1} \circ \ldots \circ \delta_{j_i}$ in degree $i$.
Via this isomorphism the differential on the right hand side
induces a differential $\text{d}$ on the left hand side.
By our sign rules we have $\text{d}(1) = - \sum f_j \delta_j$.
Since $\delta_j : \mathcal{A}^\bullet \to \mathcal{A}^\bullet[1]$
is a morphism of complexes, it follows that
$$
\text{d}(\delta_{j_1} \ldots \delta_{j_i}) =
(- \sum f_j \delta_j )\delta_{j_1} \ldots \delta_{j_i}
$$
Observe that we have $\text{d} = \sum f_j \delta_j$ on the differential
graded algebra $\mathcal{A}$. Therefore the map defined by the rule
$$
1 \circ \delta_{j_1} \ldots \delta_{j_i} \longmapsto
(\delta_{j_1} \circ \ldots \circ \delta_{j_i})(\xi_1 \ldots \xi_r)
$$
will define an isomorphism of complexes
$$
\SheafHom^\bullet(\mathcal{A}^\bullet, \mathcal{O}_Y)
\longrightarrow \mathcal{A}^\bullet[-r]
$$
if $r$ is odd and commuting with differentials up to sign if $r$ is even.
In any case these complexes have isomorphic cohomology, which shows the
desired vanishing. The isomorphism on cohomology in degree $r$
under the map
$$
\SheafHom^\bullet(\mathcal{A}^\bullet, \mathcal{O}_Y)
\longrightarrow
\SheafHom^\bullet(\mathcal{A}^\bullet, \mathcal{O}_X)
$$
also follows in a straightforward manner from this.
(We observe that our choice of conventions regarding
Koszul complexes does intervene in the definition
of the isomorphism
$R\SheafHom_{\mathcal{O}_X}(\mathcal{O}_X, \mathcal{O}_Y) =
\wedge^r \mathcal{N}[-r]$.)
\end{proof}
\end{document}
